% !Mode:: "TeX:UTF-8"

\chapter{总结与展望}\label{ch:6}
\section{总结}

随着“三高”慢性病患者数量不断增加，围绕血压、血糖、血脂管理的日常健康咨询需求持续增长。与一次性诊疗不同，“三高”管理更强调长期随访、重复咨询和生活方式干预，因此患者在日常生活中往往会提出大量细致且连续的问题。传统线下医疗服务在即时响应、高频咨询和长期陪伴方面存在一定局限，而普通互联网检索又容易出现信息来源分散、专业性不足和结论不一致等问题。与此同时，大语言模型虽然在自然语言理解与生成方面展现出较强能力，但若直接应用于医学问答场景，仍然可能出现知识幻觉、回答依据不足以及安全边界不清等风险。因此，研究如何将外部医学知识与大语言模型有效结合，构建兼顾准确性和可解释性的医学问答系统。

本文主要研究内容和结论如下：
\begin{enumerate}[label=(\arabic*)]
    \item 针对医学问答系统对知识可靠性和可追溯性的要求，本文完成了面向“三高”场景的医学知识库构建工作。研究中对政策文件、临床指南、专家共识及相关科普资料进行了采集、清洗、格式统一和结构整理，并结合医学文本特点设计了语义切分方法，在提高知识片段可检索性的同时尽量保留上下文语义完整性。此外，本文还组织了文档来源、标题层级和时间信息等元数据，为后续检索、引用展示和知识追溯提供了支撑。该部分工作为整个RAG问答系统提供了可靠的外部知识基础。
    \item 针对医学问题中同时存在术语精确匹配需求和语义理解需求的特点，本文设计了结合BM25与向量检索的混合检索方法，并引入重排机制对候选证据进行进一步筛选与排序。实验结果表明，混合检索结合重排策略在召回效果上优于单一BM25检索和单一向量检索，能够更有效地兼顾医学术语精确性与语义相关性，为后续答案生成阶段提供质量更高的检索证据。
    \item 针对基础大语言模型直接应用于医学问答时存在的回答风格不稳定、事实表达不够严谨和安全边界把握不足等问题，本文采用LoRA参数高效微调和DPO偏好优化相结合的方法进行领域适配。其中，LoRA用于以较低训练成本实现模型对医学问答任务的基础适应，DPO则进一步通过偏好数据优化模型对高质量回答的倾向性。评测结果表明，经过DPO优化后的模型在综合得分、事实正确性、完整性和安全性等维度上整体优于基础模型，说明偏好优化能够有效提升医学问答场景下的回答质量与稳健性。
    \item 在系统实现方面，本文完成了从前端页面到后端服务、从模型推理到知识维护的整体工程设计。系统前端围绕登录注册、问答主页和知识库管理页面展开，实现了用户身份验证、多轮对话交互、流式回答展示和引用证据可视化等功能；后端服务承担会话管理、消息持久化、模型调用编排和知识库运维接口等任务；模型服务则负责将前文设计的检索方法、Prompt模板和生成模型推理过程组织为统一的在线问答流程。通过这种分层架构设计，系统实现了知识检索、答案生成、证据返回和多轮交互等功能。
\end{enumerate}

\section{展望}

尽管本文已经完成系统研究与实现，并在知识库构建、混合检索、模型优化和系统落地等方面取得了一定成果，但仍有若干问题值得在后续工作中进一步完善。

第一，模型优化方面，当前偏好训练数据的规模和覆盖范围仍然有限，对复杂医学问题、模糊表达以及高风险边界场景的约束能力还有进一步提升空间。后续可继续扩充高质量偏好数据，引入更加细粒度的优劣标注方式，并探索监督微调、偏好优化与检索增强的联合训练路径，从而进一步提高模型在医学问答场景中的稳定性、鲁棒性和泛化能力。

第二，知识库与检索机制方面，当前系统主要围绕“三高”慢性病构建知识资源和检索流程，知识范围仍相对集中。后续可在保证资料权威性和可追溯性的前提下，进一步扩展至更多慢性病或其他医学专科场景，同时完善知识库的增量更新、版本管理和索引自动重建机制。此外，还可以继续研究更细粒度的文本切分、更高质量的重排策略以及更有效的证据筛选方法，以进一步提升检索结果与最终回答之间的协同效果。

第三，系统应用与功能扩展方面，当前系统已经支持文本问答、结构化健康指标补充和检测报告上传等功能，但在更贴近真实医疗咨询场景的个体化服务方面仍有拓展空间。后续可进一步结合用户历史对话、长期健康记录和多模态输入信息，探索更加个性化的健康管理问答方式；同时，也可在前端交互、医生辅助审核和安全提示机制等方面继续完善，使系统在实际应用中不仅能够“回答问题”，还能够更好地服务于慢性病随访、健康管理和风险提示等场景。
